\section{Proof Sketches}

This section provides informal proof sketches supporting the propositions stated in the preceding section. The arguments are intentionally non-exhaustive and schematic. Their purpose is to indicate why the propositions follow from the formal framing and the definition of STOP, without introducing full proofs or additional formal machinery.

\subsection*{Sketch for Proposition 1 (Model-Relativity of STOP)}
The definition of STOP is formulated relative to a specified enterprise model. Since admissibility is evaluated only within the constraints and state space defined by that model, any statement about the existence or non-existence of continuation necessarily inherits this relativity. A change of model alters the admissible continuation space and may therefore alter the diagnostic outcome.

\subsection*{Sketch for Proposition 2 (STOP as Admissibility Exhaustion)}
By definition, continuation is analytically valid only if at least one admissible continuation state exists. STOP is defined as the negation of this condition. The equivalence follows directly: if the admissible continuation space is empty, continuation is excluded; if it is non-empty, STOP does not occur. This argument relies on the notion of constraint-bounded state spaces, in which admissibility is exhausted once no permissible states remain \cite{ashby1956}.

\subsection*{Sketch for Proposition 3 (Independence from Analyst Competence)}
The analytical process evaluates admissibility within a fixed model. If no admissible continuation exists, further analytical effort cannot generate one without altering the model itself. The occurrence of STOP therefore reflects the structure of the model rather than the quality or quantity of analytical work applied to it.

\subsection*{Sketch for Proposition 4 (Non-Equivalence of STOP and Inaction)}
STOP is a statement about representability within a model, not about real-world behavior. The distinction follows from the separation between model-level diagnostics and organizational action. A STOP outcome constrains what can be asserted analytically but does not logically entail any particular response outside the model. This separation aligns with established distinctions between inquiry and action in systems analysis \cite{churchman1971}.

\subsection*{Sketch for Proposition 5 (Symmetry of Analytical Outcomes)}
Both STOP and continuation are determined by the same admissibility criterion. Neither outcome is assumed in advance; each arises depending on whether the admissible continuation space is empty or non-empty. This symmetry follows from treating continuation as a candidate outcome rather than as a default assumption.

\subsection*{Sketch for Proposition 6 (Model Revision Implication)}
If STOP occurs, then no admissible continuation exists under the current model. Any analytically valid continuation would therefore require a change in the modeling assumptions, constraints, or state space. Such changes constitute model revision rather than further analysis within the existing model and reflect a shift in the representational frame rather than an extension of inference \cite{rosen1985}.

\subsection*{Sketch for Proposition 7 (Diagnostic Informational Content)}
The exhaustion of admissible continuation states identifies the effective boundaries imposed by the model. This boundary information is non-trivial: it reveals which combinations of constraints cannot be jointly satisfied. STOP thus conveys information about the structure and limits of the modeled system.
